\chapter{Results and Discussion}\label{ch:results}

This chapter reports the evaluation. All network results come from ns-3. Each configuration runs with
several seeds, which change SUMO's trips, the choice of attackers and the radio's randomness. Tables
give the mean and the 95\,\% confidence interval over the seeds.

\section{Experimental Setup}

\begin{table}[htbp]\centering\small
\caption{Common parameters of the network experiments}\label{tab:setup}
\resizebox{\textwidth}{!}{%
\begin{tabular}{@{}lp{8.5cm}@{}}
\toprule
Parameter & Value \\
\midrule
Simulated time per run & 120\,s (180\,s for the compromised RSUs, 60\,s on LuST) \\
Scenarios & grid with 50, 100, 150 vehicles; LuST at 06:00, 14:00, 08:00 \\
RSUs & 9 (grid, 450\,m apart), 300\,m nominal range \\
Seeds & 5 (grid comparison), 10 (compromised RSUs), 3 (others) \\
Attackers & 10\,\% of the vehicles, random position offsets from 30\,s \\
Reports to revoke ($K$), reveal threshold & 2 RSUs; 5 of 9 organizations \\
Ticket lifetime, pseudonyms per vehicle & 30\,s; 5 (20 with mix zones) \\
Ledger in sweeps & in-process, 500\,ms commit (Section~\ref{sec:res-ledger}) \\
Processors & RSU at measured speed, vehicle 5$\times$ slower \\
Vehicle verification & BSMs of senders within 100\,m \\
\bottomrule
\end{tabular}}
\end{table}

Table~\ref{tab:setup} lists the common parameters. The compared variants are the Phase~I scheme, its
optimised profiles (Table~\ref{tab:profiles}) and the two baselines of Section~\ref{sec:eval-framework}.

\textbf{Metrics.}
\begin{itemize}
  \item \emph{Ticket}: time from sending a login (or renewal) to holding a valid ticket. \emph{Ticket
  after 10\,s} leaves out the start of a run.
  \item \emph{Safety traffic}: BSM and TESLA-key bytes per vehicle per second.
  \item \emph{PDR}: share of broadcasts received by the nodes within 300\,m of the sender.
  \item \emph{CBR}: channel busy ratio measured at the RSUs.
  \item \emph{V2V accept}: time from sending a BSM to a neighbouring vehicle accepting it.
  \item \emph{Revoked accepted}: share of BSMs from revoked senders that vehicles still accepted.
  \item \emph{Caught}: share of attackers reported to the ledger.
  \item \emph{Honest revoked}: honest vehicles that lost a pseudonym to revocation.
  \item \emph{Tracked}: share of pseudonym changes the eavesdropper links correctly.
\end{itemize}

\section{Comparison on the Grid}\label{sec:res-grid}

\gentab{tab_compare}
\gentab{tab_compare_sec}

Tables~\ref{tab:compare} and~\ref{tab:compare-sec} compare all variants at \R{cmp}{vehicles}
vehicles, and Figure~\ref{fig:grid} shows the trends over vehicle density.

\textbf{Latency.} The Phase~I scheme issues a ticket in \R{cmp/Phase-I}{ticket_ms_mean}\,ms,
v3-fast in \R{cmp/v3-fast}{ticket_ms_mean}\,ms (\R{cmp/derived}{ticket_ms_mean_cut_pct}\,\% less),
and the ECDSA-PKI baseline in \R{cmp/ECDSA-PKI}{ticket_ms_mean}\,ms. The zero-knowledge login is
therefore as fast as presenting a certificate. Centralising the decision (Central-TA) costs
\R{cmp/Central-TA}{ticket_ms_mean}\,ms, because every login crosses the backbone and queues at one
authority. A vehicle accepts a neighbour's BSM after \R{cmp/v3-fast}{v2v_accept_ms_mean}\,ms with
v3-fast against \R{cmp/Phase-I}{v2v_accept_ms_mean}\,ms with Phase-I: smaller messages spend less
time on the air, and the vehicle processor has less to verify.

\textbf{Channel load.} This is the largest effect. Phase-I's \R{cmp/Phase-I}{bsm_bytes}-byte BSMs at
10\,Hz put \R{cmp/Phase-I}{safety_bytes_per_vehicle_s}\,B/s per vehicle on the channel. v3-fast sends
\R{cmp/v3-fast}{safety_bytes_per_vehicle_s}\,B/s, \R{cmp/derived}{safety_factor} times less, by
combining smaller messages with dead-reckoning triggering. The channel busy ratio falls from
\R{cmp/Phase-I}{cbr_rsu_mean} to \R{cmp/v3-fast}{cbr_rsu_mean}, and the delivery ratio rises from
\R{cmp/Phase-I}{pdr} to \R{cmp/v3-fast}{pdr} because fewer frames collide.

\textbf{Revocation.} Without information from the ledger, vehicles in Phase-I accepted every BSM from
revoked senders (\R{cmp/Phase-I}{v2v_revoked_accepted_ratio}). With the revocation filter in RSU
beacons, v3-fast vehicles accepted \R{cmp/v3-fast}{v2v_revoked_accepted_ratio}. Every variant caught
every attacker, and no honest vehicle was revoked.

\textbf{Privacy.} Without mix zones, the eavesdropper links \R{cmp/v3-fast}{tracking_success} of the
pseudonym changes, because a pseudonym change alone does not hide a vehicle's trajectory. v3-private
lowers this to \R{cmp/v3-private}{tracking_success}, at the price of a ticket time of
\R{cmp/v3-private}{ticket_ms_mean}\,ms, the cost of the BBS proof.

\gentab{tab_significance}

Table~\ref{tab:significance} tests the differences: a bold entry means that the difference between
the challenger and Phase-I is significant at the 5\,\% level (Welch's t-test over the seeds).

\begin{figure}[htbp]\centering
\genfig[\textwidth]{grid_overview}
\caption{Grid scenario over vehicle density (mean and 95\,\% CI over 5 seeds)}\label{fig:grid}
\end{figure}

\gentab{tab_compare_all}

\section{Real Traffic: LuST}\label{sec:res-lust}

\gentab{tab_lust}

\begin{figure}[htbp]\centering
\genfig[\textwidth]{lust_overview}
\caption{LuST traffic at three times of day (mean and 95\,\% CI over 3 seeds)}\label{fig:lust}
\end{figure}

Table~\ref{tab:lust} and Figure~\ref{fig:lust} repeat the comparison of Phase-I and v3-fast on real
Luxembourg traffic, with 60\,s runs, since a run of the 08:00 peak with over 800 vehicles takes more
than an hour to simulate. Two differences from the grid
matter. First, a cut of real traffic starts with hundreds of vehicles already on the road, which all
log in during the first seconds. The RSUs' processors keep up (a login costs them well under a
millisecond), but the radio does not: hundreds of login requests of about 680 bytes, each with its
Merkle path and proof, and the RSUs' ticket responses compete for the channel and wait in the 802.11p
queues. In a 15\,s run of the 08:00 peak, a login request spent about 110\,ms and a ticket response
about 240\,ms on the air during the first three seconds. The mean ticket time therefore includes
this start-up burst. The ticket time after 10\,s shows the steady state: in the
08:00 peak, v3-fast issues tickets in \R{lust:peak/v3-fast}{ticket_ms_steady_mean}\,ms once the
burst has passed. Second, real roads with many lanes and junctions give more manoeuvres, so
dead-reckoning triggering saves less than on the grid. At the peak, v3-fast still puts
\R{lust:peak/v3-fast}{safety_bytes_per_vehicle_s}\,B/s per vehicle on the channel against
\R{lust:peak/Phase-I}{safety_bytes_per_vehicle_s}\,B/s for Phase-I, with channel busy ratios of
\R{lust:peak/v3-fast}{cbr_rsu_mean} and \R{lust:peak/Phase-I}{cbr_rsu_mean}.

Two results differ from the grid in v3-fast's disfavour. At the peak, its channel busy ratio is
above the 0.30 at which adaptive signing switches to TESLA, so a neighbour can accept a BSM only once
its key is disclosed: V2V acceptance takes \R{lust:peak/v3-fast}{v2v_accept_ms_mean}\,ms against
\R{lust:peak/Phase-I}{v2v_accept_ms_mean}\,ms for Phase-I. This is the price of a channel that still
delivers \R{lust:peak/v3-fast}{pdr} of the broadcasts, where Phase-I delivers
\R{lust:peak/Phase-I}{pdr}. Second, between 85 and 89\,\% of the attackers are caught on LuST,
for both variants, against all of them on the grid: the attack starts at 30\,s of a 60\,s run, and
an attacker that leaves the coverage of the RSUs before two of them have reported it is not revoked
in the run. No honest vehicle was revoked.

\section{Ablation: What Each Mechanism Contributes}\label{sec:res-ablation}

\gentab{tab_ablation}
\gentab{tab_ablation_sec}

Tables~\ref{tab:ablation} and~\ref{tab:ablation-sec} start from v2 and add one mechanism at a time. It separates the
mechanisms that reduce bytes on the air (TESLA, implicit certificates, dead reckoning) from those
that reduce latency or handovers (clusters, renewal, batch verification under load), and those that
change security properties (revocation filter, signal-strength check) or privacy (BBS, mix zones).
For example, TESLA alone brings the safety traffic from \R{abl/v2}{safety_bytes_per_vehicle_s} to
\R{abl/v2 +tesla}{safety_bytes_per_vehicle_s}\,B/s, dead reckoning alone to
\R{abl/v2 +dead-reckoning}{safety_bytes_per_vehicle_s}\,B/s, and the revocation filter alone lowers
the share of BSMs from revoked senders that vehicles accept from
\R{abl/v2}{v2v_revoked_accepted_ratio} to \R{abl/v2 +revocation-filter}{v2v_revoked_accepted_ratio}.
The learned detector in hybrid mode is listed for completeness. It is not part of the profiles: a model
trained on one road network flagged honest vehicles on another, whereas the rules and the
signal-strength check transfer across cities.

\section{Misbehaviour Detection}\label{sec:res-detect}

\gentab{tab_veremi}

On VeReMi Extension (\R{veremi}{messages} messages from \R{veremi}{vehicles} vehicles,
Table~\ref{tab:veremi}), the plausibility rules alone reach an F1 score of \R{veremi}{rules_f1} with
no false positives. They miss subtle attacks such as small constant offsets, which stay within
physical limits. Learned models over the same features reach \R{veremi}{rf_f1} (random forest~\cite{breiman2001rf}), and
gradient boosting with a history window \R{veremi}{hgb_window_f1} at a false-positive rate of
\R{veremi}{hgb_window_fpr}.

\gentab{tab_sensitivity_attack}
\gentab{tab_sensitivity_attackers}

In the simulation, Tables~\ref{tab:sensitivity-attack} and~\ref{tab:sensitivity-attackers} show
detection per attack type and for an increasing share of attackers. Every attack type is caught,
and no honest vehicle is revoked, up to 30\,\% attackers. The two-report rule matters here: a
single RSU's false alarm has no consequence, so honest vehicles keep their pseudonyms even when one
RSU misjudges them.

\section{Compromised RSUs}\label{sec:res-trust}

\gentab{tab_trust}

Table~\ref{tab:trust} lets a quarter of the RSUs report random honest pseudonyms every 5\,s. With
plain $K = 2$ counting, two compromised RSUs together can revoke an honest pseudonym. With
trust-weighted reports, an RSU's weight falls once investigations show its reports were wrong,
after which its reports no longer suffice to revoke anyone. Its first false reports still count,
so weighting does not bring the number of honest vehicles revoked to zero: it falls from 5.8 to
2.4 per run for v2, while every attacker is still caught.

\section{Privacy}\label{sec:res-privacy}

\gentab{tab_mixzone_tradeoff}

\begin{figure}[htbp]\centering
\genfig[\textwidth]{mixzone}
\caption{Mix zones: longest silent period (s) against tracking and silence}\label{fig:mixzone}
\end{figure}

Pseudonym changes alone do not stop a global eavesdropper. In the grid comparison, the tracker links
\R{cmp/v3-fast}{tracking_success} of the changes made without mix zones. Table~\ref{tab:mixzone-tradeoff}
varies the longest silent period of the cooperative mix zones. Longer silence makes linking harder,
but a silent vehicle is invisible to its neighbours, so the table also reports the share of time
vehicles were silent and the 95th percentile of the position error neighbours see. This is the
safety--privacy trade-off that a deployment has to set.

\section{Blockchain Performance}\label{sec:res-ledger}

\gentab{tab_ledger}

Table~\ref{tab:ledger} measures the real Fabric~3.1.5 network under a steady load of
\texttt{RecordAuth} transactions, the most frequent write. With Raft and Fabric's default 500\,ms
batch timeout, the median commit takes \R{ledger}{raft_500ms_5_p50}\,ms, dominated by waiting for the
block to be cut. With a 50\,ms timeout, Raft commits in \R{ledger}{raft_50ms_5_p50}\,ms and SmartBFT
in \R{ledger}{bft_50ms_5_p50}\,ms at 5 transactions per second. At 20 per second, Byzantine fault
tolerance costs more: SmartBFT's median rises to \R{ledger}{bft_50ms_20_p50}\,ms against
\R{ledger}{raft_50ms_20_p50}\,ms for Raft. No transaction failed. These latencies confirm the design
principle: a commit takes tens to hundreds of milliseconds, too slow for the path of a safety message,
but fast enough for authorization records and revocations, which the RSUs bridge with PREAUTH notes
and quarantine.

\section{Sensitivity to the Vehicle Processor}\label{sec:res-sens}

\gentab{tab_sensitivity_obu}

The speed of real on-board units varies widely: software implementations on embedded processors
verify tens of signatures per second, while units with a cryptographic accelerator verify thousands.
Instead of relying on a single assumed factor, Table~\ref{tab:sensitivity-obu} repeats the comparison
with vehicle processors as fast as the development PC, 5 times slower, and 10 times slower.

\section{Emergency Vehicles}\label{sec:res-emergency}

\gentab{tab_emergency}
\gentab{tab_emergency_sec}

\begin{figure}[htbp]\centering
\genfig[\textwidth]{emergency}
\caption{Emergency vehicles on LuST: ticket times and warning coverage}\label{fig:emergency}
\end{figure}

Tables~\ref{tab:emergency} and~\ref{tab:emergency-sec} evaluate the extension of Section~\ref{sec:emergency} on the 14:00 LuST
traffic (v3-fast, 60\,s runs), where
the start-up burst creates long radio queues at the RSUs. The first compares the ticket time of emergency vehicles with
that of all other vehicles, with and without the priority lane. The second reports the share of vehicles
within 300\,m of an emergency vehicle that held a verified alert from it at most 1\,s old, the alert
delay, and the outcome of the fake-emergency attack. The design requirement is that no fake alert is
ever accepted, and that every fake emergency login is denied by the policy or rejected because the
pseudonym is already quarantined.

\section{Comparison with Published Schemes}\label{sec:res-literature}

Published schemes report their results on different simulators, hardware and libraries, so their
numbers cannot be placed next to ours directly: BPA~\cite{li2024bpa}, for example, measured an
elliptic-curve point multiplication at 0.233\,ms, our implementation at 0.016\,ms. The usual fair
comparison counts the cryptographic operations of one authentication and prices them with one set
of unit costs. Table~\ref{tab:literature} does this with BPA's published unit costs, the only set
that covers all the schemes. For BPA, ZAMA~\cite{xi2022zama} and B-TSCA~\cite{wang2021btsca} it
copies the operation counts, times and message sizes that BPA publishes (its Tables 4 and 5). For
our scheme it counts the operations of the protocol in Chapter~\ref{ch:design} and prices only the
point multiplications and hashes, which dominate; the message sizes are measured in the ns-3 runs.

\begin{table}[htbp]\centering\small\setlength{\tabcolsep}{4pt}
\caption{Cost of one authentication and one re-authentication (handover), priced with BPA's unit
costs (point multiplication 0.233\,ms, hash 0.0055\,ms, pairing 4.756\,ms)}\label{tab:literature}
\resizebox{\ifdim\width>\textwidth\textwidth\else\width\fi}{!}{%
\begin{tabular}{lllrllrr}
\toprule
Scheme & Ledger on the & \multicolumn{2}{l}{Authentication} & \multicolumn{2}{l}{Re-authentication} &
\multicolumn{2}{l}{Messages (B)} \\
\cmidrule(lr){3-4}\cmidrule(lr){5-6}\cmidrule(lr){7-8}
 & login path & operations & ms & operations & ms & auth. & re-auth. \\
\midrule
ZAMA~\cite{xi2022zama} & no & $11E + 5M + \dots$ & 2.64 & $2E + \dots$ & 1.76 & 1200 & 472 \\
B-TSCA~\cite{wang2021btsca} & lookup & $6E + 2P + 4H + \dots$ & 10.11 & $6E + 2P + 2H + \dots$ & 10.10 & 72 & 136 \\
BPA~\cite{li2024bpa} & re-auth. & $8\mathit{PM} + 3H + \dots$ & 1.91 & $7\mathit{PM} + H + \dots$ & 1.65 & 386 & 240 \\
This work & no & $8\mathit{PM} + 15H$ & 1.95 & $10\mathit{PM} + 6H$ & 2.36 & 716 & 410 \\
\bottomrule
\end{tabular}}
\par\smallskip\parbox{\textwidth}{\scriptsize $\mathit{PM}$: point multiplication, $H$: hash, $E$: modular
exponentiation, $M$: modular multiplication, $P$: bilinear pairing; ``$\dots$'' are the cheaper
operations listed in~\cite{li2024bpa}. Rows 1--3 and their times are as published in~\cite{li2024bpa}
(its Table~5 gives 386\,B for BPA's authentication, its Figure~8 368\,B). This work: authentication
= vehicle checks the beacon (2\,PM) and proves (1\,PM), RSU checks the Merkle path (10\,H) and the
proof (2\,PM) and signs the ticket (1\,PM), vehicle checks the ticket (2\,PM); handover = vehicle
checks the beacon and signs the request (3\,PM), the new RSU checks the old ticket and the request
(4\,PM) and signs a new ticket (1\,PM), vehicle checks it (2\,PM). Message sizes: request plus
ticket response, measured in ns-3.}
\end{table}

The comparison is mixed, and Table~\ref{tab:literature} shows why. In computation our login is on a
par with BPA, which also builds on elliptic-curve operations (eight point multiplications each),
cheaper than ZAMA, whose zero-knowledge proof uses modular exponentiations, and about five times
cheaper than B-TSCA, whose two pairings alone cost 9.5\,ms. Our messages are larger: the Merkle path
(nine hashes, 288\,B) travels with every proof, which is what lets the RSU check membership without
learning which registered vehicle it is talking to. Our handover costs more than BPA's
re-authentication, because the new RSU checks a signed ticket and a signed request instead of
looking up a shared trust key; RSU clusters (Section~\ref{sec:res-ablation}) avoid most of these
handovers instead.

The largest difference is architectural. BAAIoV~\cite{alsaeed2025baaiov} authenticates every vehicle
through a Hyperledger Fabric transaction and reports 150\,ms at 50 vehicles rising to 187\,ms at
150 vehicles, measured with SUMO and Fabric but without a radio model. Our design keeps the ledger
off the login path: on the grid with 150 vehicles a login completes in
\R{cmp/v3-fast}{ticket_ms_mean}\,ms end to end over 802.11p, while the ledger records it in the
background in the tens to hundreds of milliseconds that Table~\ref{tab:ledger} measures.
The two blockchain-and-ZKP schemes closest to ours, Wang et~al.~\cite{wang2025privacy} and the
emergency-vehicle scheme of Li et~al.~\cite{li2026emergency}, are compared in Chapter~\ref{ch:lit}
by their properties only.

\section{Discussion}

The research question of the Phase~I review asked whether ZKP authentication combined with
permissioned-blockchain authorization can work without unacceptable latency, computation or
communication overhead. On the measured network the answer is yes, under two conditions. First, the
ledger must stay off the safety-message path. Proofs are verified at the RSU in a fraction of a
millisecond, while ledger commits take tens to hundreds of milliseconds and are bridged by
PREAUTH and quarantine. Second, per-message authentication must be cheap. The Phase~I design was
dominated by its 330-byte BSMs; the optimised profile is \R{cmp/derived}{safety_factor} times lighter
on the channel and faster than the certificate baseline, without weakening the security model.

\textbf{Threats to validity.}
\begin{itemize}
  \item Radio and traffic are simulated. ns-3 and SUMO are the accepted tools for this purpose, and
  LuST adds real traffic, but no field test was done.
  \item Processing costs were measured on a PC and scaled for vehicles. The processor sensitivity
  sweep bounds the effect.
  \item The parameter sweeps use an in-process ledger with the measured Raft commit time, and the
  real Fabric network was measured separately on one machine. A wide-area deployment adds network
  latency to commits.
  \item The investigation that drives trust weights is simulated with ground truth.
  \item Attackers follow the VeReMi attack types and a fake-emergency attack. Adaptive attackers who
  stay just inside the plausibility limits are harder to detect, as the VeReMi results for small
  offsets show.
\end{itemize}
